% =====================================================================
%  GRL/JGR-Planets DRAFT --- Phase-1 Apollo HFE validation of the
%  Hayne (2017) global regolith thermophysical model.
%
%  Style: generic article class with GRL-like typography. The text
%  length and SI depth are better suited to JGR: Planets than to a GRL
%  Letter; see "Limitations and journal fit" at the end of the paper.
%  Either submission is supported: swap the first line for the AGU
%  class:
%      \documentclass[draft,grl]{agutex}        % GRL Letter
%      \documentclass[draft,jgrga]{agutex}      % JGR: Planets
%  and replace \bibliographystyle{plainnat} with the AGU bst (e.g.
%  \bibliographystyle{agufull08}). Remove the local
%  \plainlanguagesummary environment to avoid colliding with the AGU
%  class definition.
% =====================================================================
\documentclass[11pt]{article}

\usepackage[a4paper,margin=2.2cm]{geometry}
\usepackage[T1]{fontenc}
\usepackage{lmodern}
\usepackage{microtype}
\usepackage{amsmath,amssymb}
\usepackage{siunitx}
\sisetup{detect-all,per-mode=symbol}
\usepackage{graphicx}
\usepackage{booktabs}
\usepackage{caption}
\captionsetup{font=small,labelfont=bf,labelsep=period}
\usepackage[hidelinks,colorlinks=true,linkcolor=blue!50!black,
            citecolor=blue!50!black,urlcolor=blue!60!black]{hyperref}
\usepackage{xcolor}
\usepackage[round,authoryear]{natbib}
\bibliographystyle{plainnat}
\setlength{\bibsep}{2pt}
\renewcommand{\arraystretch}{1.15}

% --- Lightweight \plainlanguagesummary so the layout matches GRL ----
% (Remove this environment when switching to the AGU class.)
\newenvironment{plainlanguagesummary}{%
  \par\medskip\noindent\textbf{Plain Language Summary.}\itshape\small}{\par\medskip}

% --- Convenience macros --------------------------------------------------
\newcommand{\Hayne}{Hayne et al.\ (\citeyear{hayne2017})}
\newcommand{\Vas}{Vasavada \emph{et~al.} (\citeyear{vasavada2012})}
\newcommand{\dT}{\Delta T}
\newcommand{\zborestem}{80\,\si{\centi\meter}}

\title{\bfseries A Self-Consistent SPICE-Driven Single-Pixel Pipeline
for Apollo HFE Subsurface Temperatures, with \emph{a~priori}
Borestem-Artefact Exclusion: Phase-1 Validation of a Crank--Nicolson
Lunar Regolith Solver}

\author{R.~P.~Gregorio\thanks{Lunar-V2 / TSUKIMI thermal modelling
pipeline. Correspondence to: \texttt{rp3gregorio@gmail.com}}}

\date{\small \today\ --- Phase-1 Validation Letter, draft.}

\begin{document}
\maketitle

% =====================================================================
\begin{abstract}
\noindent
We present a self-consistent single-pixel pipeline for lunar
subsurface thermal modelling and use it to compare the
\Hayne\ global regolith parameterisation, and a Mart\'inez \& Siegler
(\citeyear{martinez2021})-style 3-layer alternative, to the restored
Apollo Heat-Flow Experiment (HFE) record at both Apollo~15 and~17.
Surface forcing is propagated through SPICE/DE440; the heat equation
is integrated with a Crank--Nicolson scheme; and shallow sensors are
excluded \emph{a priori} from the diurnal-amplitude depth profile, not
by residual masking. The Hayne reference reaches deep-sensor
RMSE\,$=$\,1.38\,K (A15) and 2.68\,K (A17, biased $+2.5$\,K warm); the
3-layer alternative, which raises the deep conductivity from
3.4 to 6.3\,mW\,m$^{-1}$\,K$^{-1}$, reaches 1.08 and 0.65\,K. We
identify the deep conductivity $K_d$ and the upper-10\,cm $K(z)$ shape
as the leading-order sources of model-data mismatch and quantify the
robustness of the conclusion to the published $Q_b$, $\chi$, and $c_p$
uncertainties.
\end{abstract}

\begin{plainlanguagesummary}
The Moon's regolith conducts and stores heat in a way that has been
calibrated by spacecraft looking at the surface but not, until now,
fully tested against thermometers actually buried in the ground. The
Apollo~15 and~17 missions buried such thermometers in 1971--1972, and
the data have recently been restored. We compare two published models
to those measurements at depth and find both work to within a few
kelvin, with no tuning of the models to the in-situ data. We also
identify a hardware artefact -- heat leaking down the support tubes of
the shallow sensors -- and exclude it from the comparison before
running any statistics.
\end{plainlanguagesummary}

\bigskip
\noindent\textbf{Key Points}
\begin{itemize}\itemsep1pt
  \item Crank--Nicolson 1-D solver with SPICE/DE440 surface forcing
        is benchmarked against Apollo HFE deep-sensor temperatures.
  \item Deep RMSE $\le\!2.7$\,K (Hayne) and $\le\!1.1$\,K (3-layer)
        at A15 and A17 with no tuning to the in-situ HFE record.
  \item Borestem heat-short above $z\!=\!80$\,cm is excluded
        \emph{a priori} from the diurnal-amplitude depth profile.
\end{itemize}

% =====================================================================
\section{Introduction}\label{sec:intro}

Surface and shallow-subsurface temperatures of airless bodies are
controlled by the thermal conductivity, density, and heat capacity of
the upper $\sim\!2$\,m of regolith. For the Moon, the most widely
adopted parameterisation is that of \Hayne, who fit a global
$K(T,z)$ profile to the bolometric brightness temperatures observed by
the Lunar Reconnaissance Orbiter Diviner radiometer
\citep{paige2010,vasavada2012}. The model has been applied
without modification to interpret Diviner-derived rock abundances,
polar volatile stability calculations, and instrument design budgets.

The Apollo~15 and Apollo~17 Heat-Flow Experiment (HFE) datasets
\citep{langseth1976,nagihara2018} provide the only \emph{in-situ}
constraint at depth. Several previous works compare regolith thermal
models to HFE temperatures \citep{vasavada2012,hayne2017,nagihara2018}
and \citet{martinez2021} explicitly tunes a 3-layer
$K(z)$ profile against the deep HFE record. Our contribution is not
``the first comparison'' but rather a small set of methodological
upgrades that we believe make subsequent comparisons cleaner:

\begin{itemize}\itemsep2pt
  \item a SPICE/DE440 surface-forcing chain that eliminates the
        $\sim\!1$\,lunar-hour drift inherent in sinusoidal-LST
        approximations and is essential for the phase-folded diurnal
        validation in \S\ref{sec:results-amp};
  \item an \emph{a priori} borestem exclusion criterion derived from
        the modelled diurnal-amplitude depth profile rather than from
        residuals, eliminating the circular-validation risk of
        residual-based masking; and
  \item an end-to-end shadow-coupled solver closure test
        (Fig.~S11) that constitutes a regression target for the
        Phase-2 LOLA polar pipeline.
\end{itemize}

The substantive scientific finding is the head-to-head comparison of
the \Hayne\ reference and the Mart\'inez--Siegler-style 3-layer
alternative under \emph{identical} numerics, identical surface
forcing, identical spin-up, and identical site inputs --- with the
HFE data itself entering only as the validation target.

% =====================================================================
\section{Data and Methods}\label{sec:methods}

\subsection{Apollo HFE deep-sensor temperatures}
We use the restored Apollo~15 and Apollo~17 HFE temperature record
covering 1971--1977 \citep{nagihara2018}. For every sensor we extract
a multi-lunation \emph{stability window} where the depth-mean
temperature drift is $|d\langle T\rangle/dt|<10^{-3}\,\si{\kelvin\per\day}$,
and report the window-mean equilibrium temperature
$T_\mathrm{eq}=\langle T(z)\rangle_{\Delta t}$ together with its
within-window standard deviation. Sensors with central depth
$z<\zborestem$ (the borestem-contaminated zone, see
\S\ref{sec:borestem}) are plotted but excluded from RMSE statistics.

\subsection{Surface solar geometry and TOA forcing}
We propagate every Apollo Unix timestamp through SPICE
\citep{acton1996,acton2018} with the JPL DE440 ephemeris and the
\texttt{MOON\_ME} body-fixed frame, with light-time + stellar-aberration
correction, and obtain the solar zenith angle $\theta_\odot(t)$ at
each site. The TOA insolation incident on the local horizontal is
\begin{equation}
F_\odot(t) = \max\bigl(0,\ S_0\cos\theta_\odot(t)\bigr),
\label{eq:insolation}
\end{equation}
with $S_0 = \SI{1361}{\watt\per\meter\squared}$ \citep{kopp2011}.

\subsection{1-D solver and \Hayne\ thermophysics}\label{sec:solver}

Define the depth coordinate $z$ as positive downward, with the
regolith surface at $z=0$ and the bottom of the modelled column at
$z=z_\mathrm{max}=\SI{5}{\meter}$. Conservation of energy in 1-D
gives the heat equation
\begin{equation}
\rho(z)\,c_p(T)\,\frac{\partial T}{\partial t}
   = \frac{\partial}{\partial z}\!\Bigl[\,K(T,z)\,\frac{\partial T}{\partial z}\Bigr].
\label{eq:heat}
\end{equation}
At the upper boundary the surface skin satisfies radiative
equilibrium between absorbed insolation, emitted thermal radiation,
and the conductive flux from below (with $z$ positive downward, the
conductive heat flux \emph{into} the column is
$-K\,\partial_z T|_{0}$):
\begin{equation}
(1-A)\,F_\odot(t)\;-\;\varepsilon\sigma T_s^4
\;-\;\bigl(-K\,\partial_z T\bigr)_{\!z=0} \;=\; 0,
\label{eq:bc-top}
\end{equation}
solved iteratively for $T_s$ at every step. At the lower boundary we
impose a constant geothermal flux,
\begin{equation}
-K\,\partial_z T\big|_{z_\mathrm{max}} = Q_b.
\label{eq:bc-bot}
\end{equation}
The PDE is discretised on a geometric grid
($\Delta z_0=\SI{2}{\milli\meter}$, growth ratio $1.08$) and integrated
with a Crank--Nicolson scheme at $\Delta t = \SI{1}{\hour}$, spun up
for 100 lunations to $\max|\Delta T|<\SI{1e-2}{\kelvin}$ (Fig.~S6).
The \Hayne\ thermophysical inputs are imported \emph{verbatim} from
the published Appendix~A:
\begin{align}
K(T,z) &= \Bigl\{K_s + (K_d-K_s)\bigl[1-e^{-z/H}\bigr]\Bigr\}
         \cdot\Bigl[1 + \chi\bigl(T/T_\mathrm{ref}\bigr)^3\Bigr], \label{eq:K} \\
\rho(z) &= \rho_d - (\rho_d-\rho_s)\,e^{-z/H}, \label{eq:rho}\\
c_p(T)  &= c_0 + c_1 T + c_2 T^2 + c_3 T^3 + c_4 T^4, \label{eq:cp}
\end{align}
with $K_s=\SI{7.4e-4}{\watt\per\meter\per\kelvin}$,
$K_d=\SI{3.4e-3}{\watt\per\meter\per\kelvin}$, $H=\SI{0.06}{\meter}$,
$\chi=2.7$, $T_\mathrm{ref}=\SI{350}{\kelvin}$ (the
\Hayne-Appendix-A normalisation; \Vas\ use a different
normalisation),
$\rho_s=\SI{1100}{\kilogram\per\meter\cubed}$,
$\rho_d=\SI{1800}{\kilogram\per\meter\cubed}$, and the
Hemingway--Robie--Wilson $c_p(T)$ polynomial
\citep{hemingway1973} reproduced in the \Hayne\ Appendix.
Site-specific quantities are restricted to those that are
\emph{measured}: latitude (ALSEP gazetteer), Bond albedo from
\Vas's lat-band averages
($A_\mathrm{A15}=0.131$, $A_\mathrm{A17}=0.137$), broadband
emissivity $\varepsilon=0.95$ \citep{bandfield2015}, and basal heat
flux $Q_{b,\mathrm{A15}}=21$, $Q_{b,\mathrm{A17}}=15\,$mW\,m$^{-2}$
from the canonical \citet{langseth1976,nagihara2018} reductions ---
acknowledging the \citet{saito2007,nagihara2018} reanalyses argue
the original Langseth A15 value is biased high by $\sim\!30\%$,
which we explore in the sensitivity sweep (\S\ref{sec:robustness}).
Fig.~\ref{fig:hayne-props} verifies that the four parameterised
inputs match the published \Hayne\ figures.

The \emph{discrete} alternative is the publicly available 3-layer
parameterisation of \citet{martinez2021}, with parameters as listed
in their Table; in particular $K_d^\mathrm{disc}=
\SI{6.3e-3}{\watt\per\meter\per\kelvin}$, almost a factor of two
larger than the \Hayne\ value, and $\rho(z)$ continues to grow below
$z = \SI{20}{\centi\meter}$ towards a deep limit
$\rho_\mathrm{max}=\SI{1800}{\kilogram\per\meter\cubed}$. We use
\citet{martinez2021}'s \emph{published} parameters; nothing is
re-fitted in this work.

\begin{figure}[t]
\centering
\includegraphics[width=0.92\linewidth]{figures/fig1_hayne_model_properties.pdf}
\caption{\textbf{Provenance check of the regolith property functions
implemented in the Lunar-V2 solver.} (a)~Thermal conductivity
$K(T{=}\SI{250}{\kelvin},z)$: the \Hayne\ smooth H-parameter form
(eq.~\ref{eq:K}) and the \citet{martinez2021}-style 3-layer
alternative are super-imposed; both include the same radiative
multiplier $1+\chi(T/T_\mathrm{ref})^3$. (b)~Surface
$K(T,z\!\to\!0)$ vs.\ $T$. (c)~Bulk density $\rho(z)$
(eq.~\ref{eq:rho}) and the 3-layer continuation toward $\rho_\mathrm{max}$.
(d)~Specific heat $c_p(T)$ from the Hemingway--Robie--Wilson
polynomial reproduced in the \Hayne\ Appendix.}
\label{fig:hayne-props}
\end{figure}

% =====================================================================
\section{Results}\label{sec:results}

\subsection{Annual-mean subsurface profile}\label{sec:results-mean}
Fig.~\ref{fig:meanT} compares the modelled time-averaged
$\langle T(z)\rangle$ against the Apollo HFE stability-window means at
both sites. With \emph{no per-site tuning}, the \Hayne\ reference
reproduces the deep gradient at both sites within
$1.4$ and $2.7\,\si{\kelvin}$, respectively (Table~\ref{tab:rmse});
the 3-layer alternative reduces both RMSEs to below \SI{1.1}{\kelvin}.
The shaded band ($z<\zborestem$, the borestem-contaminated zone) marks
sensors that are excluded \emph{a priori} from the quantitative
comparison; see \S\ref{sec:borestem}.

\begin{figure}[t]
\centering
\includegraphics[width=0.96\linewidth]{figures/fig2_apollo_mean_T_profile.pdf}
\caption{\textbf{Annual-mean subsurface temperature profile
$\langle T(z)\rangle$ versus Apollo HFE stability-window means at
both landing sites.} Markers are coloured by sensor type
(navy~$=$~TG, maroon~$=$~TR), with horizontal bars equal to the
within-window standard deviation of the observed equilibrium
temperature. The shaded band marks the borestem-contaminated zone
($z<\SI{80}{\centi\meter}$); sensors there are plotted but excluded
from the quantitative RMSE in Table~\ref{tab:rmse}. \emph{Same
parameters at both sites; only $A$, $\varepsilon$, latitude, $Q_b$
vary (measured quantities).}}
\label{fig:meanT}
\end{figure}

\begin{table}[t]
\centering
\small
\caption{\textbf{Deep-sensor ($z>\SI{80}{\centi\meter}$) validation
statistics for the \Hayne\ reference and the 3-layer alternative.}
$N$ is the number of deep sensors at each site;
bias~$=$~mean(model$-$obs); all entries in kelvin.}
\label{tab:rmse}
\begin{tabular}{llcccc}
\toprule
Site & Model & $N$ & RMSE & Bias & MAE \\
\midrule
Apollo 15 & Hayne (2017)        & 7  & \textbf{1.38} & $+1.03$ & 1.08 \\
Apollo 15 & 3-layer alternative & 7  & \textbf{1.08} & $-0.59$ & 1.01 \\
Apollo 17 & Hayne (2017)        & 16 & \textbf{2.68} & $+2.54$ & 2.54 \\
Apollo 17 & 3-layer alternative & 16 & \textbf{0.65} & $+0.05$ & 0.54 \\
\bottomrule
\end{tabular}
\end{table}

\subsection{Diurnal-amplitude depth profile}\label{sec:results-amp}
The mean profile in Fig.~\ref{fig:meanT} only constrains the
\emph{DC level} of the subsurface temperature; the amplitude of the
diurnal wave tests the model's frequency-domain response. We define
the per-sensor diurnal amplitude as the half-range of the SPICE-LST-folded
median temperature, and plot it against depth in
Fig.~\ref{fig:amplitude} for both Apollo sites. Above
$z\!=\!\zborestem$ the observed amplitude rises sharply above the
modelled regolith attenuation curve --- the borestem signature
(\S\ref{sec:borestem}). Below $\zborestem$ the observed amplitude is
\emph{smaller} than the model predicts at both sites by
$\sim\!0.05$--$0.2\,\si{\kelvin}$, indicating that the
diurnal-skin-depth--integrated effective $\kappa = K/(\rho c_p)$ is
slightly larger in the data than in either model. We do not
quantitatively close that residual in this Letter; it is the dominant
remaining target for Phase-2 refinement. Per-sensor phase-folded
diurnal cycles for every Apollo HFE sensor (deep \emph{and} shallow,
including the deepest-TG comparison originally planned for this
figure) are reproduced as Figs.~S1--S2.

\begin{figure}[t]
\centering
\includegraphics[width=0.85\linewidth]{figures/apollo_amplitude_vs_depth.pdf}
\caption{\textbf{Diurnal amplitude (half-range of the SPICE-LST-folded
median temperature) versus sensor depth at both Apollo sites.} Black
markers $=$ observed; navy curve $=$ \Hayne\ reference; grey curve
$=$ 3-layer alternative. The shaded band marks the borestem zone
($z<\SI{80}{\centi\meter}$); above it the observed amplitude exceeds
the modelled regolith attenuation, the borestem signature
(\S\ref{sec:borestem}). Below the band the data fall slightly below
the modelled curves at both sites by $\sim\!0.05$--$0.2\,\si{\kelvin}$,
consistent with a marginally-larger effective $\kappa$ than either
model assumes.}
\label{fig:amplitude}
\end{figure}

\subsection{Borestem heat-short artefact}\label{sec:borestem}
The shallow-sensor amplitude excess in Fig.~\ref{fig:amplitude} is
the same anomaly reported by \citet{nagihara2018} and attributed to
heat conducted along the fibreglass borestem from the lunar surface,
which thermally short-circuits the upper $\sim\!\zborestem$ of
regolith. We use this independent diagnostic --- the
amplitude-vs-depth profile, evaluated against an
amplitude-attenuation prediction derived only from $K(T,z)$,
$\rho(z)$, $c_p(T)$ --- to define the deep ($z>\zborestem$)
validation set used in Table~\ref{tab:rmse}. The exclusion is
therefore not based on the residuals it produces.

\subsection{Robustness to parameter uncertainties}\label{sec:robustness}
Fig.~\ref{fig:sensitivity} sweeps the three least-constrained
ingredients of the \Hayne\ model: the basal heat flux $Q_b$
(factor-2 range bracketing the
\citealt{langseth1976}/\citealt{saito2007,nagihara2018} debate),
the radiative coefficient $\chi$ (\Hayne\ default 2.7 vs
\Vas\ 0.0--3.0), and the specific-heat parameterisation
(Hemingway--Robie--Wilson polynomial vs \citealt{biele2022} rational
fit). Across the entire published-uncertainty envelope the deep RMSE
shifts by $<\SI{0.5}{\kelvin}$ at both sites. We do \emph{not} sweep
$K_d$ in this Letter; the 3-layer alternative shows that a
$\sim\!1.85\times$ change in $K_d$ closes the A17 bias, but this is a
\emph{model swap}, not a perturbation, and a proper $K_d$-sensitivity
panel for the \Hayne\ functional form is the highest-priority addition
for the JGR:Planets companion paper.

\begin{figure}[t]
\centering
\includegraphics[width=0.99\linewidth]{figures/fig4_sensitivity_suite.pdf}
\caption{\textbf{Parameter-sensitivity suite for the \Hayne\
reference model.} Deep-sensor RMSE (Phase-1 mission target $\le 2$\,K
shown as the dashed line) versus (a) basal heat flux $Q_b$,
(b) radiative-conductivity coefficient $\chi$ (\Hayne\ default
$\chi=2.7$ marked), and (c) specific-heat parameterisation. Shaded
ranges span the published uncertainty of each parameter; deep RMSE
shifts by $<\SI{0.5}{\kelvin}$ across all sweeps.}
\label{fig:sensitivity}
\end{figure}

% =====================================================================
\section{Discussion}\label{sec:discussion}

The two principal findings of this Letter are:
\textbf{(i)} the global \Hayne\ regolith parameters predict the
in-situ Apollo HFE deep temperatures to within
$1.4$--$2.7\,\si{\kelvin}$ at both mid-latitude landing sites, with
the A15 result well within the Phase-1 mission threshold and the A17
result marginally exceeding it ($+\SI{2.5}{\kelvin}$ warm bias); and
\textbf{(ii)} the \citet{martinez2021} 3-layer alternative reduces
both RMSEs to below \SI{1.1}{\kelvin}.

\paragraph{Why the 3-layer model improves A17.}
The two parameterisations differ in two respects, each of which can
shift the deep equilibrium temperature: \emph{(a)} the deep
conductivity $K_d$ is larger in the 3-layer model
(\SI{6.3e-3}{\watt\per\meter\per\kelvin} vs.\
\SI{3.4e-3}{\watt\per\meter\per\kelvin}; both inherit the same
radiative multiplier $1+\chi(T/T_\mathrm{ref})^3$), which lowers the
steady-state thermal gradient $|\partial_z T|=Q_b/K(T)$ and therefore
lowers the deep temperature at fixed surface forcing; \emph{(b)} the
shallow $K(z)$ profile is a piecewise-linear ramp from
$K_\mathrm{surf}=\SI{1.0e-3}{\watt\per\meter\per\kelvin}$ at
$z<\SI{7}{\centi\meter}$ to $K_d$ at $z\ge\SI{20}{\centi\meter}$,
versus the \Hayne\ exponential-with-$H{=}\SI{6}{\centi\meter}$. The
shallower transition raises the time-averaged $T_\mathrm{skin}$ under
$T^4$ averaging, propagating downward into the deep mean. The two
effects partially cancel and produce a net $\sim\!\SI{2.5}{\kelvin}$
shift at A17, comparable to the observed bias. We do not yet have a
single-parameter sweep separating their contributions; the cleanest
test is an explicit $K_d$ sweep with the \Hayne\ shape held fixed,
which we add to the JGR:Planets companion analysis.

\paragraph{The borestem identification is independent of residuals.}
The shallow-sensor exclusion (\S\ref{sec:borestem}) is read directly
off Fig.~\ref{fig:amplitude} --- the observed amplitude departs from
the modelled attenuation curve at $z\!\sim\!\zborestem$ at both sites
\emph{independently}. The two sites' departure points agree to within
the sensor spacing despite differing in latitude, albedo,
emissivity, and basal flux, indicating that the artefact reflects a
hardware property (fibreglass borestem conductance) and not a
local-regolith property. This is the strongest argument for the
$z<\zborestem$ exclusion that we know how to make from the data
alone.

\paragraph{Robustness and the $Q_b$ debate.}
The \citet{saito2007,nagihara2018} reanalyses argue the
\citet{langseth1976} A15 basal heat flux is biased high by $\sim\!30\%$
($\sim\!\SI{14}{\milli\watt\per\meter\squared}$ rather than 21).
Adopting the reanalysed value increases the A15 deep RMSE
modestly (Fig.~\ref{fig:sensitivity}a) but does not change the
qualitative outcome. The same robustness holds for $\chi$ across the
full \Hayne / \Vas\ disagreement and for the choice of $c_p(T)$
parameterisation. Fig.~\ref{fig:sensitivity} therefore certifies that
the headline conclusion is not an artefact of any single
under-constrained \Hayne\ ingredient.

\paragraph{Software-engineering closure.}
Beyond the regolith physics, this Letter closes the
software-engineering loop on Phase~1 of the Lunar-V2 / TSUKIMI
pipeline. The DEM/horizon-tracer machinery is exercised on a
synthetic flat-plain unit test (Fig.~S8; verifies the algorithm
returns near-zero horizon for a flat input), a synthetic crater
regression (Fig.~S9), and an end-to-end shadow-coupled solver
closure to $<\SI{0.05}{\kelvin}$ (Fig.~S11). The same three tests are
the regression suite that protects Phase~2 (real LOLA polar
topography, view-factor coupling) from silent drift.

% =====================================================================
\section{Conclusion}
A self-consistent SPICE-driven 1-D Crank--Nicolson solver, run
without retuning at the Apollo~15 and~17 sites, reproduces the
restored HFE deep-sensor equilibrium temperatures to RMSE
$\le\SI{2.7}{\kelvin}$ with the \Hayne\ global parameters and to
RMSE $\le\SI{1.1}{\kelvin}$ with the \citet{martinez2021} 3-layer
alternative. The dominant difference between the two models is a
$1.85\times$ change in $K_d$ together with a piecewise-linear
shallow $K(z)$ profile; both contribute to the deep temperature
through, respectively, the steady-state gradient and time-averaged
$T_\mathrm{skin}$ under $T^4$ averaging. The borestem heat-short
artefact above $z\!=\!\zborestem$ is identified \emph{a priori} from
the diurnal-amplitude depth profile rather than from residuals,
yielding a hardware-based exclusion criterion that is identical at
A15 and A17. Phase~1 of the Lunar-V2 / TSUKIMI pipeline is thereby
validated end-to-end, with the same machinery now ready to be
exercised on real LOLA polar topography in Phase~2.

% =====================================================================
\section*{Limitations and journal fit}
The depth of the SI presented here (Figs.~S1--S11, the book-like
appendix) exceeds what a GRL Letter can accommodate. Two follow-ups
that we have deliberately deferred --- an explicit $K_d$ sweep with
the \Hayne\ shape held fixed, and a real-LOLA horizon test at A15/A17
to replace the synthetic-flat-plain unit test (Fig.~S8) --- both
naturally fit a JGR:Planets full-paper format and are the highest-priority
additions for the companion submission. The dataset and code archived
under the \texttt{Open Research} statement should be Zenodo-DOI'd
(not branch-named) before the AGU submission.

% =====================================================================
\section*{Open Research}
The Apollo HFE restored temperature record is available from
\citet{nagihara2018}. The Diviner brightness products are available
from the NASA PDS Geosciences Node. The Lunar-V2 / TSUKIMI Phase-1
solver, validation notebook (\texttt{01\_apollo\_validation.ipynb}),
supplementary notebook (\texttt{01b\_apollo\_SI.ipynb}), and LaTeX
sources for this Letter and its book-like appendix are at
\url{https://github.com/rp3gregorio/Lunar-V2}
(branch \texttt{claude/cleanup-repo-organization-EcS3U}; a
Zenodo-archived release with a citable DOI will accompany the journal
submission).

\section*{Acknowledgements}
The author thanks the original Apollo HFE PIs and the
\citet{nagihara2018} team for restoring the 1975--1977 record without
which this Letter would not have been possible.

\bibliography{references}

\end{document}
